\documentclass[11pt,a4paper]{article}
\usepackage[top=2.3cm,bottom=2.3cm,left=2.2cm,right=2.2cm]{geometry}

\usepackage{amsmath,amssymb}
\usepackage{fontspec}
\setmainfont{Liberation Serif}
\setsansfont{Liberation Sans}
\setmonofont{DejaVu Sans Mono}[Scale=0.84]

\usepackage{xcolor}
\definecolor{brandblue}{RGB}{20, 55, 105}
\definecolor{accentblue}{RGB}{35, 95, 165}
\definecolor{codebg}{RGB}{248, 249, 251}
\definecolor{codeframe}{RGB}{215, 222, 230}
\definecolor{javakeyword}{RGB}{0, 45, 170}
\definecolor{javastring}{RGB}{5, 120, 20}
\definecolor{javacomment}{RGB}{120, 120, 120}
\definecolor{javanumber}{RGB}{20, 75, 220}

\definecolor{noteborder}{RGB}{30, 110, 65}
\definecolor{notebg}{RGB}{242, 249, 245}
\definecolor{warnborder}{RGB}{185, 75, 10}
\definecolor{warnbg}{RGB}{254, 248, 240}
\definecolor{tipborder}{RGB}{30, 95, 160}
\definecolor{tipbg}{RGB}{242, 247, 254}

\usepackage{fancyhdr}
\usepackage{lastpage}
\pagestyle{fancy}
\fancyhf{}
\lhead{\parbox[t]{0.58\textwidth}{\raggedright\small\sffamily\color{brandblue}\textbf{T15 --- Tipi Generici, Wildcard e Type Erasure}}}
\rhead{\small\sffamily\color{gray}Programmazione II -- UniVR (2025/2026)}
\cfoot{\small\sffamily Pagina \thepage\ di \pageref{LastPage}}
\renewcommand{\headrulewidth}{0.5pt}
\renewcommand{\headrule}{\hbox to\headwidth{\color{brandblue!70!black}\leaders\hrule height \headrulewidth\hfill}}

\usepackage{titlesec}
\titleformat{\section}{\Large\bfseries\sffamily\color{brandblue}}{\thesection}{0.8em}{}[\titlerule]
\titleformat{\subsection}{\large\bfseries\sffamily\color{accentblue}}{\thesubsection}{0.8em}{}
\titleformat{\subsubsection}{\normalsize\bfseries\sffamily\color{brandblue!85!black}}{\thesubsubsection}{0.8em}{}

\usepackage{needspace}
\usepackage{fancyvrb}
\DefineVerbatimEnvironment{asciibox}{Verbatim}{
  fontsize=\fontsize{7.5pt}{9.2pt}\selectfont,
  frame=single,
  rulecolor=\color{codeframe},
  fillcolor=\color{codebg},
  framesep=2.5mm
}
\usepackage{listings}
\lstset{
  backgroundcolor=\color{codebg},
  frame=single,
  rulecolor=\color{codeframe},
  basicstyle=\ttfamily\footnotesize,
  keywordstyle=\color{javakeyword}\bfseries,
  stringstyle=\color{javastring},
  commentstyle=\color{javacomment}\itshape,
  numberstyle=\tiny\color{gray},
  numbers=left,
  numbersep=7pt,
  stepnumber=1,
  breaklines=true,
  breakatwhitespace=true,
  showstringspaces=false,
  tabsize=4,
  xleftmargin=12pt,
  framexleftmargin=12pt
}

\newcommand{\passthrough}[1]{#1}
\providecommand{\tightlist}{\setlength{\itemsep}{0pt}\setlength{\parskip}{0pt}}

\usepackage{tcolorbox}
\tcbuselibrary{skins,breakable}
\newtcolorbox{studynote}[1][Nota]{
  enhanced,
  breakable,
  colback=notebg,
  colframe=noteborder,
  fonttitle=\bfseries\sffamily\small,
  coltitle=white,
  title={#1},
  arc=2mm,
  boxrule=0.8pt,
  left=9pt, right=9pt, top=7pt, bottom=7pt
}

\newtcolorbox{studywarning}[1][Attenzione]{
  enhanced,
  breakable,
  colback=warnbg,
  colframe=warnborder,
  fonttitle=\bfseries\sffamily\small,
  coltitle=white,
  title={#1},
  arc=2mm,
  boxrule=0.8pt,
  left=9pt, right=9pt, top=7pt, bottom=7pt
}

\newtcolorbox{studytip}[1][Suggerimento]{
  enhanced,
  breakable,
  colback=tipbg,
  colframe=tipborder,
  fonttitle=\bfseries\sffamily\small,
  coltitle=white,
  title={#1},
  arc=2mm,
  boxrule=0.8pt,
  left=9pt, right=9pt, top=7pt, bottom=7pt
}

\usepackage{booktabs}
\usepackage{longtable}
\usepackage{array}
\usepackage{calc}

\usepackage{hyperref}
\hypersetup{
  colorlinks=true,
  linkcolor=accentblue,
  urlcolor=accentblue,
  citecolor=brandblue,
  pdfborder={0 0 0}
}

\title{\Huge\bfseries\sffamily\color{brandblue} T15 --- Tipi Generici, Wildcard e Type Erasure \\[0.3em] \large\sffamily Polimorfismo Parametrico, Bounded Type Parameters, Invarianza, Wildcard (PECS) e Type Erasure}
\author{\textbf{Docente:} Prof. Michele Pasqua \quad---\quad \textbf{Appunti Revisione:} Wally}
\date{A.A. 2025/2026 -- Universit\`a degli Studi di Verona}

\begin{document}
\maketitle
\thispagestyle{fancy}
\vspace{-1.2em}
\noindent\textcolor{brandblue}{\rule{\textwidth}{0.8pt}}
\vspace{1.2em}

\begin{center}\rule{0.5\linewidth}{0.5pt}\end{center}

\subsection{1. Analisi Teorica Approfondita (Slide
T15)}\label{analisi-teorica-approfondita-slide-t15}

La lezione 15 introduce formalmente i \textbf{Tipi Generici}
(\emph{Generics}), introdotti in Java 5 per abilitare il
\textbf{polimorfismo parametrico}, eliminare la fragilità dei cast a
\passthrough{\lstinline!Object!} e garantire la \textbf{Type Safety} a
tempo di compilazione.

\begin{center}\rule{0.5\linewidth}{0.5pt}\end{center}

\subsubsection{\texorpdfstring{1.1 L'Idea Guida \& La Tazza
\texttt{Cup\textless{}T\textgreater{}} (Slide
2-3)}{1.1 L'Idea Guida \& La Tazza Cup\textless T\textgreater{} (Slide 2-3)}}\label{lidea-guida-la-tazza-cupt-slide-2-3}

Nella Slide 3 campeggia l'immagine di una tazza con la scritta:
\[\mathbf{Cup\langle T\rangle}\] accompagnata dal motto: \textbf{``Same
container, different content''}. Un contenitore (una tazza, una coppia,
una lista, un albero) possiede una logica strutturale identica a
prescindere dal tipo di dato che custodisce (caffè, tè, stringhe,
interi, date). Prima di Java 5, per ottenere contenitori riusabili si
usava \passthrough{\lstinline!Object!}, con conseguenze disastrose
sull'affidabilità del software.

\begin{center}\rule{0.5\linewidth}{0.5pt}\end{center}

\subsubsection{\texorpdfstring{1.2 Polimorfismo Parametrico vs Approccio
Basato su \texttt{Object} (Slide
4-7)}{1.2 Polimorfismo Parametrico vs Approccio Basato su Object (Slide 4-7)}}\label{polimorfismo-parametrico-vs-approccio-basato-su-object-slide-4-7}

Consideriamo la funzione identità matematica \(\lambda x.\, x\): * Se
riceve un \passthrough{\lstinline!int!}, restituisce lo stesso
\passthrough{\lstinline!int!}. * Se riceve una
\passthrough{\lstinline!String!}, restituisce la stessa
\passthrough{\lstinline!String!}.

\paragraph{\texorpdfstring{L'approccio ingenuo (pre-Java 5): segnaposto
\texttt{Object}}{L'approccio ingenuo (pre-Java 5): segnaposto Object}}\label{lapproccio-ingenuo-pre-java-5-segnaposto-object}

\begin{lstlisting}[language=Java]
public class Pair {
    private Object first, second;
    public Pair(Object first, Object second) {
        this.first = first;
        this.second = second;
    }
    public Object getFirst() { return first; }
    public Object getSecond() { return second; }
}
\end{lstlisting}

\begin{itemize}
\item
  \textbf{Casting obbligatorio}:

\begin{lstlisting}[language=Java]
Pair p = new Pair("one", "two");
String first = (String) p.getFirst(); // Cast esplicito verboso e tedioso
\end{lstlisting}
\item
  \textbf{Il disastro della mancanza di Type Safety}:

\begin{lstlisting}[language=Java]
Pair q = new Pair(1, 2); // Autoboxing da int a Integer
String second = (String) q.getSecond(); // COMPILA SENZA ERRORI!
// Ma a RUNTIME crasha con: java.lang.ClassCastException: class java.lang.Integer cannot be cast to class java.lang.String
\end{lstlisting}
\end{itemize}

\begin{center}\rule{0.5\linewidth}{0.5pt}\end{center}

\subsubsection{1.3 Classi e Metodi Generici (Slide
8-10)}\label{classi-e-metodi-generici-slide-8-10}

Con i Generics, la classe introduce una o più \textbf{variabili di tipo}
(\emph{Type Parameters}):

\begin{lstlisting}[language=Java]
public class Pair<T> {
    private T first, second;
    public Pair(T first, T second) {
        this.first = first;
        this.second = second;
    }
    public T getFirst() { return first; }
    public T getSecond() { return second; }
}
\end{lstlisting}

\begin{itemize}
\item
  \textbf{Sintassi d'uso}:

\begin{lstlisting}[language=Java]
Pair<String> p = new Pair<String>("one", "two");
Pair<Integer> q = new Pair<Integer>(1, 2);
\end{lstlisting}
\item
  \textbf{Controllo a compile-time}:

\begin{lstlisting}[language=Java]
String first = p.getFirst(); // Nessun cast necessario: p restituisce String garantita!
String second = q.getSecond(); // COMPILE-TIME ERROR: Type mismatch (Integer non convertibile a String)
\end{lstlisting}

  L'errore viene intercettato prima ancora di mandare in esecuzione il
  programma!
\item
  \textbf{Regola fondamentale sui Tipi Primitivi}: \textbf{Un parametro
  di tipo \(T\) può essere sostituito solo da tipi riferimento
  (oggetti), MAI da tipi primitivi!} Non è consentito scrivere
  \passthrough{\lstinline!Pair<int>!}, ma si deve usare il wrapper
  \passthrough{\lstinline!Pair<Integer>!}.
\item
  \textbf{Convenzioni di Naming}: Si usano lettere maiuscole singole:

  \begin{itemize}
  \tightlist
  \item
    \passthrough{\lstinline!T!}: \emph{Type} generico.
  \item
    \passthrough{\lstinline!E!}: \emph{Element} (usato nelle collezioni:
    \passthrough{\lstinline!List<E>!},
    \passthrough{\lstinline!Set<E>!}).
  \item
    \passthrough{\lstinline!K, V!}: \emph{Key} e \emph{Value} (usato
    nelle mappe: \passthrough{\lstinline!Map<K, V>!}).
  \item
    \passthrough{\lstinline!N!}: \emph{Number}.
  \item
    \passthrough{\lstinline!S, U!}: Tipi secondari ausiliari.
  \end{itemize}
\end{itemize}

\begin{center}\rule{0.5\linewidth}{0.5pt}\end{center}

\subsubsection{1.4 Interfacce della Libreria Standard Riscritte con i
Generics (Slide
11-13)}\label{interfacce-della-libreria-standard-riscritte-con-i-generics-slide-11-13}

\begin{itemize}
\item
  \textbf{\passthrough{\lstinline!Comparable<T>!}}:

\begin{lstlisting}[language=Java]
public interface Comparable<T> {
    int compareTo(T obj);
}
\end{lstlisting}

  Consente alla classe che la implementa (es.
  \passthrough{\lstinline!Student implements Comparable<Student>!}) di
  ricevere direttamente un parametro di tipo
  \passthrough{\lstinline!Student!} nel metodo
  \passthrough{\lstinline!compareTo(Student other)!}, eliminando
  controlli \passthrough{\lstinline!instanceof!} e downcast forzati.
\item
  \textbf{\passthrough{\lstinline!Iterable<E>!} e
  \passthrough{\lstinline!Iterator<E>!}}:

\begin{lstlisting}[language=Java]
public interface Iterable<E> {
    Iterator<E> iterator();
}
public interface Iterator<E> {
    boolean hasNext();
    E next();
}
\end{lstlisting}

  Permette a costrutti come il ciclo \passthrough{\lstinline!for-each!}
  di iterare in modo fortemente tipizzato senza alcun cast:

\begin{lstlisting}[language=Java]
for (Student student : classRoom) {
    System.out.println(student.toString());
}
\end{lstlisting}
\end{itemize}

\begin{center}\rule{0.5\linewidth}{0.5pt}\end{center}

\subsubsection{\texorpdfstring{1.5 Il Diamond Operator
\texttt{\textless{}\textgreater{}} (Slide
14)}{1.5 Il Diamond Operator \textless\textgreater{} (Slide 14)}}\label{il-diamond-operator-slide-14}

Introdotto in Java 7, evita la ridondanza tra la dichiarazione del tipo
di riferimento e l'istanziazione:

\begin{lstlisting}[language=Java]
// Java 5 e 6:
Pair<Integer> intPair = new Pair<Integer>(2, 4);

// Java 7+:
Pair<Integer> intPair = new Pair<>(2, 4); // Il compilatore inferisce il tipo Integer dai membri a sinistra
\end{lstlisting}

\begin{center}\rule{0.5\linewidth}{0.5pt}\end{center}

\subsubsection{1.6 Bounded Type Parameters: Vincolare i Parametri di
Tipo (Slide
15-17)}\label{bounded-type-parameters-vincolare-i-parametri-di-tipo-slide-15-17}

\begin{itemize}
\item
  \textbf{Unbounded Type Parameter (Slide 15)}: Se dichiariamo
  \passthrough{\lstinline!public class Pair<T>!}, \(T\) è privo di
  vincoli. Il compilatore applica il meccanismo del \textbf{Type
  Erasure} assumendo che \(T\) sia un generico
  \passthrough{\lstinline!Object!}. Di conseguenza, su un oggetto di
  tipo \(T\) è possibile invocare \textbf{esclusivamente i metodi
  definiti in \passthrough{\lstinline!Object!}}
  (\passthrough{\lstinline!equals!}, \passthrough{\lstinline!hashCode!},
  \passthrough{\lstinline!toString!},
  \passthrough{\lstinline!getClass!}). Se tentiamo di chiamare
  \passthrough{\lstinline!first.getName()!}, il compilatore rigetta il
  codice con errore:
  \passthrough{\lstinline!cannot find symbol: method getName() in T!}.
\item
  \textbf{Upper Bound con \passthrough{\lstinline!extends!} (Slide
  16-17)}: Possiamo imporre che \(T\) sia un sottotipo di una
  determinata classe o implementi una o più interfacce:

\begin{lstlisting}[language=Java]
public class Pair<T extends Person> {
    private T first, second;
    // ...
    public String getFirstName() {
        return first.getName(); // PERFETTAMENTE LEGALE: T è garantito essere almeno una Person!
    }
}
\end{lstlisting}

  \begin{itemize}
  \tightlist
  \item
    \textbf{Vincoli multipli}:
    \passthrough{\lstinline!<T extends ClassA \& InterfaceB \& InterfaceC>!}.
    La classe (se presente) deve comparire per prima, seguita dalle
    interfacce separate da \passthrough{\lstinline!\&!}.
  \item
    \emph{Nota terminologica della Slide 16}: La slide cita
    \passthrough{\lstinline!class Clazz<T super S>!}, ma in Java a
    livello di dichiarazione di classe generica è ammesso solo
    \passthrough{\lstinline!extends!}. La parola chiave
    \passthrough{\lstinline!super!} è utilizzabile
    \textbf{esclusivamente nelle Wildcards} (che vediamo subito sotto).
  \end{itemize}
\end{itemize}

\begin{center}\rule{0.5\linewidth}{0.5pt}\end{center}

\subsubsection{1.7 Ereditarietà e Generics: La Trappola dell'Invarianza
(Slide
18-20)}\label{ereditarietuxe0-e-generics-la-trappola-dellinvarianza-slide-18-20}

Questo è uno dei punti più critici dell'intero corso e delle domande
d'esame.

\needspace{7cm}
\begin{asciibox}
   Ereditarietà Normale             Ereditarietà con Generics
      ┌────────────┐                     ┌───────────────┐
      │   Person   │                     │ Pair<Person>  │
      └─────▲──────┘                     └───────▲───────┘
            │ extends                            │ (Nessun legame
            │                                    │  di parentela!)
      ┌─────┴──────┐                     ┌───────┴───────┐
      │  Student   │                     │ Pair<Student> │
      └────────────┘                     └───────────────┘
   Student È UN Person             Pair<Student> NON È Pair<Person>
\end{asciibox}

\begin{itemize}
\item
  \textbf{Invarianza dei Generics (Slide 19)}: Anche se
  \passthrough{\lstinline!Student!} è sottotipo di
  \passthrough{\lstinline!Person!},
  \textbf{\passthrough{\lstinline!Pair<Student>!} NON è un sottotipo di
  \passthrough{\lstinline!Pair<Person>!}}!

\begin{lstlisting}[language=Java]
Student s = new Student("Sam", 123);
Person p = s; // OK: principio di sostituzione di Liskov

Pair<Student> pairS = new Pair<>(s, new Student("Paul", 456));
Pair<Person> pairP = pairS; // COMPILE-TIME ERROR: Incompatible types!
\end{lstlisting}
\item
  \textbf{Perché Java vieta questo assegnamento?}: Se
  \passthrough{\lstinline!pairP = pairS!} fosse consentito, potremmo
  eseguire:

\begin{lstlisting}[language=Java]
pairP.setFirst(new Professor("Mario", 999)); // Legale su Pair<Person>, poiché un Professor è una Person!
\end{lstlisting}

  Ma poiché \passthrough{\lstinline!pairP!} e
  \passthrough{\lstinline!pairS!} punterebbero allo stesso identico
  oggetto nell'Heap, \passthrough{\lstinline!pairS!} (che pretende di
  contenere solo \passthrough{\lstinline!Student!}) si troverebbe a
  contenere un \passthrough{\lstinline!Professor!}, distruggendo la
  garanzia di tipo a runtime!
\item
  \textbf{Il Contrasto con gli Array (Array Covariance, Slide 20)}: In
  Java gli array sono storicamente \textbf{covarianti}:

\begin{lstlisting}[language=Java]
Student[] arrS = { new Student("Sam", 123) };
Person[] arrP = arrS; // COMPILA PERFETTAMENTE!
arrP[0] = new Professor("Mario", 999); // BOOM! A runtime genera: java.lang.ArrayStoreException
\end{lstlisting}

  La covarianza degli array è universalmente riconosciuta come una
  debolezza storica di Java (introdotta in Java 1.0 prima dell'avvento
  dei generics). I Generics sono stati resi \textbf{invarianti} proprio
  per risolvere questa falla a tempo di compilazione.
\end{itemize}

\begin{center}\rule{0.5\linewidth}{0.5pt}\end{center}

\subsubsection{\texorpdfstring{1.8 Wildcards:
\texttt{\textless{}?\textgreater{}},
\texttt{\textless{}?\ extends\ T\textgreater{}},
\texttt{\textless{}?\ super\ T\textgreater{}} (Slide
21-23)}{1.8 Wildcards: \textless?\textgreater, \textless? extends T\textgreater, \textless? super T\textgreater{} (Slide 21-23)}}\label{wildcards-extends-t-super-t-slide-21-23}

Per recuperare la flessibilità polimorfica senza rinunciare alla Type
Safety, Java ha introdotto il carattere jolly (\emph{Wildcard})
\passthrough{\lstinline!?!}.

\begin{enumerate}
\def\labelenumi{\arabic{enumi}.}
\tightlist
\item
  \textbf{Unbounded Wildcard (\passthrough{\lstinline!<?>!}, Slide 21)}:

  \begin{itemize}
  \item
    Rappresenta un tipo totalmente sconosciuto:

\begin{lstlisting}[language=Java]
Pair<?> pairUnknown = pairS; // Legale!
\end{lstlisting}
  \item
    Possiamo \textbf{leggere} elementi, ma il loro tipo a compile-time è
    unicamente \passthrough{\lstinline!Object!}.
  \item
    \textbf{Non possiamo scrivere o aggiungere nulla} (tranne il
    letterale \passthrough{\lstinline!null!}), perché non sappiamo quale
    sia il tipo effettivo accettabile.
  \end{itemize}
\item
  \textbf{Upper-Bounded Wildcard
  (\passthrough{\lstinline!<? extends T>!}, Slide 22-23)}:

  \begin{itemize}
  \item
    Accetta \passthrough{\lstinline!T!} o qualsiasi sottoclasse di
    \passthrough{\lstinline!T!}:

\begin{lstlisting}[language=Java]
Pair<? extends Person> pairP = pairS; // Legale! Student extends Person
\end{lstlisting}
  \item
    Ora possiamo invocare in sicurezza tutti i metodi di
    \passthrough{\lstinline!Person!}:

\begin{lstlisting}[language=Java]
System.out.println(pairP.getFirst().getName()); // OK!
\end{lstlisting}
  \item
    \textbf{Sintassi pulita nei metodi}: Invece di dichiarare un
    parametro formale di tipo esplicito a livello di metodo:

\begin{lstlisting}[language=Java]
<T extends Person> void printPair(Pair<T> p) { ... }
\end{lstlisting}

    possiamo scrivere in modo molto più compatto e leggibile:

\begin{lstlisting}[language=Java]
void printPair(Pair<? extends Person> p) { ... }
\end{lstlisting}
  \end{itemize}
\item
  \textbf{Il Principio Guida PECS (\emph{Producer Extends, Consumer
  Super})}:

  \begin{itemize}
  \tightlist
  \item
    Se la struttura generica funge da \textbf{produttore} di dati
    (vogliamo solo \emph{leggere} dati di tipo
    \passthrough{\lstinline!T!}), si usa
    \passthrough{\lstinline!<? extends T>!}.
  \item
    Se la struttura funge da \textbf{consumatore} di dati (vogliamo
    \emph{scrivere/aggiungere} elementi di tipo
    \passthrough{\lstinline!T!}), si usa
    \passthrough{\lstinline!<? super T>!}.
  \end{itemize}
\end{enumerate}

\begin{center}\rule{0.5\linewidth}{0.5pt}\end{center}

\subsection{\texorpdfstring{2. Analisi Dettagliata del Codice
(\texttt{Lezione15/MavenDate})}{2. Analisi Dettagliata del Codice (Lezione15/MavenDate)}}\label{analisi-dettagliata-del-codice-lezione15mavendate}

Tra la Lezione 14 e la Lezione 15, il prof. Pasqua compie due modifiche
sostanziali: 1. Implementa \passthrough{\lstinline!hashCode()!} in
\href{file:///home/wally/Documenti/GitHub/Programmazione-II/25-26/Teoria-e-Esercitazioni/Codice-20260902/Lezione15/MavenDate/src/main/java/it/oop/core/Date.java\#L102-L105}{Date.java}.
2. Riscrive ed espande
\href{file:///home/wally/Documenti/GitHub/Programmazione-II/25-26/Teoria-e-Esercitazioni/Codice-20260902/Lezione15/MavenDate/src/main/java/it/oop/ui/MainDate.java}{MainDate.java}
per integrare la libreria delle \textbf{Java Collections}
(\passthrough{\lstinline!List!}, \passthrough{\lstinline!Set!},
\passthrough{\lstinline!Map!}), il contratto
\passthrough{\lstinline!equals!}/\passthrough{\lstinline!hashCode!}, e i
meccanismi di ordinamento con \passthrough{\lstinline!Comparable!} e
\passthrough{\lstinline!Comparator!} (usando classi anonime e lambdas).

\begin{center}\rule{0.5\linewidth}{0.5pt}\end{center}

\subsubsection{\texorpdfstring{2.1 Il Contratto \texttt{equals} /
\texttt{hashCode} in
\href{file:///home/wally/Documenti/GitHub/Programmazione-II/25-26/Teoria-e-Esercitazioni/Codice-20260902/Lezione15/MavenDate/src/main/java/it/oop/core/Date.java\#L102-L105}{Date.java}}{2.1 Il Contratto equals / hashCode in Date.java}}\label{il-contratto-equals-hashcode-in-date.java}

Nel passaggio da Lezione 14 a Lezione 15 viene aggiunto a
\passthrough{\lstinline!Date!}:

\begin{lstlisting}[language=Java]
    @Override
    public int hashCode() {
        return (day + month) * year;
    }
\end{lstlisting}

\paragraph{Perché questo metodo è vitale e non
opzionale?}\label{perchuxe9-questo-metodo-uxe8-vitale-e-non-opzionale}

In Java vige la regola ferrea del contratto di
\passthrough{\lstinline!Object!}: \textgreater{} \textbf{Se due oggetti
sono considerati uguali dal metodo
\passthrough{\lstinline!equals(Object)!}
(\passthrough{\lstinline!a.equals(b) == true!}), allora le chiamate a
\passthrough{\lstinline!hashCode()!} su entrambi gli oggetti DEVONO
restituire lo stesso valore intero
(\passthrough{\lstinline!a.hashCode() == b.hashCode()!}).}

Se si fa l'override di \passthrough{\lstinline!equals!} senza fare
l'override di \passthrough{\lstinline!hashCode!}, strutture dati basate
su hash come \passthrough{\lstinline!HashSet!} e
\passthrough{\lstinline!HashMap!} falliscono miseramente: considerano
due date identiche come memorizzate in \emph{bucket} differenti,
inserendo duplicati o non trovando chiavi presenti.

Nel codice di \passthrough{\lstinline!MainDate.java!}:

\begin{lstlisting}[language=Java]
AmericanDate usD1 = new AmericanDate(7, 11, 2025);
AmericanDate usD2 = new AmericanDate(7, 11, 2025);
System.out.println("usD1 ?= usD2: " + usD1.equals(usD2)); // true
System.out.println("hash(usD1): " + usD1.hashCode());      // 36450
System.out.println("hash(usD2): " + usD2.hashCode());      // 36450
\end{lstlisting}

Calcolo: \((7 + 11) \times 2025 = 18 \times 2025 = 36450\). Entrambi gli
oggetti hanno identico hash code!

\begin{center}\rule{0.5\linewidth}{0.5pt}\end{center}

\subsubsection{\texorpdfstring{2.2 \texttt{List} vs \texttt{Set}:
Accettazione vs Rifiuto dei Duplicati in
\href{file:///home/wally/Documenti/GitHub/Programmazione-II/25-26/Teoria-e-Esercitazioni/Codice-20260902/Lezione15/MavenDate/src/main/java/it/oop/ui/MainDate.java\#L43-L55}{MainDate.java}}{2.2 List vs Set: Accettazione vs Rifiuto dei Duplicati in MainDate.java}}\label{list-vs-set-accettazione-vs-rifiuto-dei-duplicati-in-maindate.java}

\begin{lstlisting}[language=Java]
// 1. Creazione di una List (ammette elementi duplicati e mantiene l'ordine di inserimento)
List<Date> dateList = new ArrayList<>();
dateList.add(new AmericanDate(7, 11, 2025));
dateList.add(new Date(9, 11, 2025));
dateList.add(new TimeStamp(0, 0, 0, 8, 11, 2025));
dateList.add(new AmericanDate(7, 11, 2025)); // Duplicato!
System.out.println("list(" + dateList.size() + "): " + dateList.toString());

// 2. Creazione di un Set (insieme matematico: nessun duplicato consentito)
Set<Date> dateSet = new HashSet<>();
dateSet.add(new AmericanDate(7, 11, 2025));
dateSet.add(new Date(9, 11, 2025));
dateSet.add(new TimeStamp(0, 0, 0, 8, 11, 2025));
dateSet.add(new AmericanDate(7, 11, 2025)); // Duplicato: viene scartato silenziosamente!
System.out.println("set(" + dateSet.size() + "): " + dateSet.toString());
\end{lstlisting}

\begin{itemize}
\tightlist
\item
  \textbf{Risultato a runtime}:

  \begin{itemize}
  \tightlist
  \item
    \passthrough{\lstinline!dateList.size()!} vale \textbf{4} (contiene
    entrambe le istanze di \passthrough{\lstinline!7/11/2025!}).
  \item
    \passthrough{\lstinline!dateSet.size()!} vale \textbf{3} (la seconda
    istanza è stata rifiutata perché \passthrough{\lstinline!equals!} ha
    restituito \passthrough{\lstinline!true!} e
    l'\passthrough{\lstinline!hashCode!} coincideva).
  \end{itemize}
\end{itemize}

\begin{center}\rule{0.5\linewidth}{0.5pt}\end{center}

\subsubsection{\texorpdfstring{2.3
\texttt{Map\textless{}K,\ V\textgreater{}} e Lambdas Registrate
Dinamicamente in
\href{file:///home/wally/Documenti/GitHub/Programmazione-II/25-26/Teoria-e-Esercitazioni/Codice-20260902/Lezione15/MavenDate/src/main/java/it/oop/ui/MainDate.java\#L61-L69}{MainDate.java}}{2.3 Map\textless K, V\textgreater{} e Lambdas Registrate Dinamicamente in MainDate.java}}\label{mapk-v-e-lambdas-registrate-dinamicamente-in-maindate.java}

\begin{lstlisting}[language=Java]
Map<String, FormattedDateConverter> converter = new HashMap<>();
// Registrazione di strategie di conversione tramite espressioni lambda
converter.put("usToIt", d -> new ItalianDate(d.getDay(), d.getMonth(), d.getYear()));
converter.put("itToUs", d -> new AmericanDate(d.getDay(), d.getMonth(), d.getYear()));

FormattedDate usDate = new AmericanDate(1, 1, 1970);
FormattedDate itDate = converter.get("usToIt").convert(usDate);
System.out.println(itDate.prettyPrint()); // Stampa "1 gennaio 1970"

for (String k : converter.keySet())
    System.out.println(converter.get(k).toString());
\end{lstlisting}

\begin{itemize}
\tightlist
\item
  \textbf{Finezza di esame sul \passthrough{\lstinline!toString()!}
  delle Lambdas}: Quando si invoca \passthrough{\lstinline!toString()!}
  su un oggetto lambda in Java, la JVM non stampa il codice sorgente,
  bensì il descrittore sintetico della classe generata a runtime
  (\passthrough{\lstinline!invokedynamic!}):
  \passthrough{\lstinline!it.oop.ui.MainDate$$Lambda/0x...!} (la
  notazione interna con i due dollari \passthrough{\lstinline!$$Lambda!}
  è la convenzione standard con cui la JVM battezza le classi sintetiche
  generate dinamicamente).
\end{itemize}

\begin{center}\rule{0.5\linewidth}{0.5pt}\end{center}

\subsubsection{\texorpdfstring{2.4 Ordinamento: \texttt{Comparable} vs
\texttt{Comparator} (Slide 11 \&
\href{file:///home/wally/Documenti/GitHub/Programmazione-II/25-26/Teoria-e-Esercitazioni/Codice-20260902/Lezione15/MavenDate/src/main/java/it/oop/ui/MainDate.java\#L70-L88}{MainDate.java})}{2.4 Ordinamento: Comparable vs Comparator (Slide 11 \& MainDate.java)}}\label{ordinamento-comparable-vs-comparator-slide-11-maindate.java}

In Java esistono due modi per ordinare oggetti: 1. \textbf{Ordinamento
Naturale con \passthrough{\lstinline!Comparable<T>!}}: * Definito
all'interno della classe stessa tramite
\passthrough{\lstinline!compareTo(T other)!}. *
\passthrough{\lstinline!Date!} implementa
\passthrough{\lstinline!Comparable<Date>!}:
\passthrough{\lstinline"java      Collections.sort(dateList); // Funziona direttamente!"}
La lista disordinata
\passthrough{\lstinline![11/7/2025, y2025m11d9, y2025m11d8, 11/7/2025]!}
viene riordinata cronologicamente:
\passthrough{\lstinline![11/7/2025, 11/7/2025, y2025m11d8[00:00:00], y2025m11d9]!}.

\begin{enumerate}
\def\labelenumi{\arabic{enumi}.}
\setcounter{enumi}{1}
\tightlist
\item
  \textbf{Ordinamento Personalizzato con
  \passthrough{\lstinline!Comparator<T>!}}:

  \begin{itemize}
  \item
    Consideriamo \passthrough{\lstinline!DateInterval!}: eredita da
    \passthrough{\lstinline!OrderedPair<Date>!}, ma \textbf{NON
    implementa \passthrough{\lstinline!Comparable<DateInterval>!}}!
  \item
    Se si prova a chiamare:

\begin{lstlisting}[language=Java]
// Collections.sort(intvList); // COMPILE-TIME ERROR:
// no suitable method found for sort(List<DateInterval>)
// reason: cannot infer type-variable(s) T (DateInterval is not Comparable)
\end{lstlisting}
  \item
    Per risolvere il problema, il prof. Pasqua applica due stili per
    fornire un \passthrough{\lstinline!Comparator<DateInterval>!}
    esterno:

    \begin{itemize}
    \item
      \textbf{Stile 1: Classe Anonima} (ordinamento per data di inizio
      \passthrough{\lstinline!left!} crescente):

\begin{lstlisting}[language=Java]
Collections.sort(intvList, new Comparator<DateInterval>() {
    @Override
    public int compare(DateInterval di1, DateInterval di2) {
        return di1.getLeft().compareTo(di2.getLeft());
    }
});
\end{lstlisting}
    \item
      \textbf{Stile 2: Espressione Lambda} (ordinamento per data di fine
      \passthrough{\lstinline!right!} decrescente):

\begin{lstlisting}[language=Java]
Collections.sort(intvList, (di1, di2) -> di2.getRight().compareTo(di1.getRight()));
\end{lstlisting}
    \end{itemize}
  \end{itemize}
\end{enumerate}

\begin{center}\rule{0.5\linewidth}{0.5pt}\end{center}

\subsection{3. Risultato di Compilazione ed Esecuzione
Reale}\label{risultato-di-compilazione-ed-esecuzione-reale}

Comando eseguito nel progetto
\href{file:///home/wally/Documenti/GitHub/Programmazione-II/25-26/Teoria-e-Esercitazioni/Codice-20260902/Lezione15/MavenDate}{Lezione15/MavenDate}:

\begin{lstlisting}[language=bash]
mvn compile exec:java -Dexec.mainClass="it.oop.ui.MainDate"
\end{lstlisting}

Output ottenuto:

\begin{lstlisting}
y2025m9d10
y2025m1d1
00:00:00
true
OOP exam on 2 febbraio 2026
date pair: [y2025m2d1 .. y2025m1d1]
Pair not ordered!
OOP exam on 2 febbraio 2026
list(4): [11/7/2025, y2025m11d9, y2025m11d8[00:00:00], 11/7/2025]
set(3): [11/7/2025, y2025m11d9, y2025m11d8[00:00:00]]
usD1 ?= usD2: true
hash(usD1): 36450
hash(usD2): 36450
1 gennaio 1970
it.oop.ui.MainDate$$Lambda/0x000000008f383b48@21043ee9
it.oop.ui.MainDate$$Lambda/0x000000008f383920@6dc5f73f
[11/7/2025, 11/7/2025, y2025m11d8[00:00:00], y2025m11d9]
[[y2025m1d1 .. y2025m1d20], [y2025m1d10 .. y2025m1d31]]
[[y2025m1d10 .. y2025m1d31], [y2025m1d1 .. y2025m1d20]]
\end{lstlisting}

\subsubsection{Decodifica puntuale
dell'output:}\label{decodifica-puntuale-delloutput}

\begin{enumerate}
\def\labelenumi{\arabic{enumi}.}
\tightlist
\item
  \passthrough{\lstinline!list(4): [...]!} vs
  \passthrough{\lstinline!set(3): [...]!} \(\rightarrow\) Dimostrazione
  pratica che \passthrough{\lstinline!HashSet!} elimina il duplicato
  \passthrough{\lstinline!11/7/2025!} grazie al corretto funzionamento
  congiunto di \passthrough{\lstinline!equals()!} e
  \passthrough{\lstinline!hashCode()!}.
\item
  \passthrough{\lstinline!hash(usD1): 36450!} e
  \passthrough{\lstinline!hash(usD2): 36450!} \(\rightarrow\) Rispetto
  formale del contratto di \passthrough{\lstinline!Object!}.
\item
  \passthrough{\lstinline!1 gennaio 1970!} \(\rightarrow\) La mappa
  \passthrough{\lstinline!converter!} recupera il converter
  \passthrough{\lstinline!"usToIt"!} memorizzato come lambda e trasforma
  \passthrough{\lstinline!1/1/1970!} in formato italiano.
\item
  \passthrough{\lstinline!it.oop.ui.MainDate$$Lambda/...!}
  \(\rightarrow\) Nome sintetico dell'oggetto funzionale generato dalla
  JVM a runtime (in cui il doppio dollaro
  \passthrough{\lstinline!$$Lambda!} rappresenta il discriminatore
  standard del bytecode).
\item
  \passthrough{\lstinline![11/7/2025, 11/7/2025, y2025m11d8[00:00:00], y2025m11d9]!}
  \(\rightarrow\) Lista ordinata secondo l'ordinamento naturale di
  \passthrough{\lstinline!Date!} (7 nov, poi 8 nov, poi 9 nov).
\item
  \passthrough{\lstinline![[y2025m1d1 .. y2025m1d20], [y2025m1d10 .. y2025m1d31]]!}
  \(\rightarrow\) Intervalli ordinati per data d'inizio crescente (1 gen
  prima del 10 gen) tramite \passthrough{\lstinline!Comparator!} in
  classe anonima.
\item
  \passthrough{\lstinline![[y2025m1d10 .. y2025m1d31], [y2025m1d1 .. y2025m1d20]]!}
  \(\rightarrow\) Intervalli ordinati per data di fine decrescente (31
  gen prima del 20 gen) tramite \passthrough{\lstinline!Comparator!} in
  espressione lambda.
\end{enumerate}

\begin{center}\rule{0.5\linewidth}{0.5pt}\end{center}

Dimmi \textbf{``vai''} per procedere con il \textbf{Blocco 15} (Lezione
16 / Slide T16 --- \emph{Collections Framework}, gerarchia completa di
\passthrough{\lstinline!List!}, \passthrough{\lstinline!Set!},
\passthrough{\lstinline!Queue!}, \passthrough{\lstinline!Map!},
complessità asintotiche, iteratori e diff in
\passthrough{\lstinline!Lezione16/MavenDate!}).


\end{document}
